\documentclass[12pt]{article}

\usepackage[letterpaper, margin=2.2cm]{geometry}
\usepackage[spanish]{babel}
\usepackage[utf8]{inputenc}
\usepackage{amsmath}
\usepackage{setspace}
\usepackage{url}

\onehalfspacing

\begin{document}

\section{KPIs medibles según ISO/IEC 25023}

Con base en el estándar ISO/IEC 25023, se definieron indicadores clave de desempeño (KPIs) para evaluar atributos de calidad relevantes en el proyecto Cal.com. Estos indicadores permiten medir objetivamente el comportamiento del sistema y verificar el cumplimiento de estándares de calidad asociados a disponibilidad, rendimiento, mantenibilidad, compatibilidad y usabilidad \cite{iso25023}. Además, la selección de estas métricas se relaciona con prácticas de medición y aseguramiento de calidad descritas por Pressman y Maxim \cite{pressman2020}.

\subsection{Disponibilidad del sistema}

\textbf{Característica ISO 25010 relacionada:} Fiabilidad

\textbf{KPI:} Disponibilidad del sistema

\textbf{Fórmula:}

\[
Disponibilidad = \frac{Tiempo\ Operativo}{Tiempo\ Operativo + Tiempo\ de\ Inactividad} \times 100
\]

\textbf{Umbral:} $\geq 99.5\%$

\textbf{Herramienta de medición:} Grafana y logs del servidor.

\subsection{Tiempo promedio de respuesta}

\textbf{Característica ISO 25010 relacionada:} Eficiencia de desempeño

\textbf{KPI:} Tiempo promedio de respuesta de APIs

\textbf{Fórmula:}

\[
TPR = \frac{\sum Tiempo\ de\ respuesta}{Numero\ de\ solicitudes}
\]

\textbf{Umbral:} $\leq 2\ segundos$

\textbf{Herramienta de medición:} Postman y Lighthouse.

\subsection{Cobertura de pruebas}

\textbf{Característica ISO 25010 relacionada:} Mantenibilidad

\textbf{KPI:} Cobertura de pruebas automatizadas

\textbf{Fórmula:}

\[
Cobertura = \frac{Lineas\ probadas}{Total\ de\ lineas\ de\ codigo} \times 100
\]

\textbf{Umbral:} $\geq 80\%$

\textbf{Herramienta de medición:} Jest y SonarQube.

\subsection{Compatibilidad con navegadores}

\textbf{Característica ISO 25010 relacionada:} Compatibilidad

\textbf{KPI:} Compatibilidad multiplataforma

\textbf{Fórmula:}

\[
Compatibilidad = \frac{Pruebas\ exitosas}{Pruebas\ totales} \times 100
\]

\textbf{Umbral:} $\geq 95\%$

\textbf{Herramienta de medición:} BrowserStack.

\subsection{Tasa de errores críticos}

\textbf{Característica ISO 25010 relacionada:} Fiabilidad

\textbf{KPI:} Errores críticos por semana

\textbf{Fórmula:}

\[
TEC = \frac{Errores\ criticos}{Semana}
\]

\textbf{Umbral:} $\leq 3$

\textbf{Herramienta de medición:} GitHub Issues y Sentry.

\subsection{Accesibilidad de la interfaz}

\textbf{Característica ISO 25010 relacionada:} Usabilidad

\textbf{KPI:} Nivel de cumplimiento WCAG

\textbf{Fórmula:}

\[
Accesibilidad = \frac{Elementos\ accesibles}{Elementos\ evaluados} \times 100
\]

\textbf{Umbral:} $\geq 90\%$

\textbf{Herramienta de medición:} Lighthouse y WAVE.

\newpage

\section{Aportes individuales}

\textbf{Diego Calva Alvarez}

Mi aporte principal en este entregable fue el desarrollo de la sección 4.3, correspondiente a los KPIs medibles según ISO/IEC 25023. Para esta parte definí indicadores específicos aplicables al proyecto Cal.com, considerando atributos de calidad como disponibilidad, rendimiento, mantenibilidad, compatibilidad, fiabilidad y usabilidad. También establecí fórmulas, umbrales numéricos y herramientas de medición para cada KPI, con el objetivo de que las métricas pudieran evaluarse de forma clara y objetiva dentro del proyecto.

\newpage

\nocite{ieee730}

\bibliographystyle{apalike}
\bibliography{referencias}

\end{document}